% Figure: drug-eluting stent cross-section with concentration profile.
% Vol VI Ch 12 sec 12.4.
\begin{figure}[htbp]
\centering
\begin{tikzpicture}[
  layer/.style={draw=engblack, minimum height=0.5cm, align=center,
    font=\scriptsize},
  arrow/.style={->, thick, engblack},
]

% Stent strut cross-section: left half is the strut, right side is tissue
% Strut (metal)
\fill[engblack!30] (-3.5, -0.6) rectangle (0, 0.6);
\node[font=\small, white] at (-1.7, 0) {\textbf{stent strut (Co-Cr)}};

% Drug-loaded polymer coating
\fill[engblue!40] (0, -0.6) rectangle (0.6, 0.6);
\node[font=\tiny, anchor=center, rotate=90] at (0.3, 0)
  {polymer + drug};

% Tissue (arterial intima / media)
\fill[engred!10] (0.6, -0.6) rectangle (5.5, 0.6);
\node[font=\small] at (3, 0) {\textbf{arterial tissue}};

% Boundary labels
\draw[engblack, thick] (0, -0.6) -- (0, 0.6);
\draw[engblack, thick] (0.6, -0.6) -- (0.6, 0.6);

% Annotation: thickness L
\draw[<->] (0, -0.9) -- (0.6, -0.9);
\node[font=\scriptsize, anchor=north] at (0.3, -0.9)
  {$L \approx 5\text{--}15\,\mu\text{m}$};

% Concentration profile on top
\draw[->, thick] (-3.5, 1.5) -- (5.7, 1.5);
\node[font=\scriptsize, anchor=west] at (5.7, 1.5) {$x$};

\draw[->, thick] (-3.5, 1.5) -- (-3.5, 4.0);
\node[font=\scriptsize, anchor=east] at (-3.5, 4.0) {$C(x, t)$};

% Concentration profile: high in coating, decreasing into tissue
\draw[engblue, very thick] plot[domain=0:0.6, samples=10]
  (\x, {3.5 - 0.5 * \x});
\draw[engblue, very thick] plot[domain=0.6:5.5, samples=40]
  (\x, {3.2 * exp(-0.7 * (\x - 0.6)) + 1.5});

% Time progression: dashed for later time
\draw[engblue, very thick, dashed] plot[domain=0:0.6, samples=10]
  (\x, {2.4 - 0.3 * \x});
\draw[engblue, very thick, dashed] plot[domain=0.6:5.5, samples=40]
  (\x, {2.0 * exp(-0.4 * (\x - 0.6)) + 1.55});

% Labels for the curves
\node[font=\scriptsize, anchor=west, engblue] at (1.0, 3.5)
  {$t_1$ (early)};
\node[font=\scriptsize, anchor=west, engblue] at (1.5, 2.4)
  {$t_2$ (late)};

% Vertical dashed lines at boundaries
\draw[engblack!50, dashed] (0, 1.5) -- (0, 4.0);
\draw[engblack!50, dashed] (0.6, 1.5) -- (0.6, 4.0);

% Bottom: cumulative release plot
\draw[->, thick] (-3.5, -2.0) -- (5.5, -2.0);
\node[font=\scriptsize, anchor=west] at (5.5, -2.0) {$t$};
\draw[->, thick] (-3.5, -2.0) -- (-3.5, -0.7);
\node[font=\scriptsize, anchor=east] at (-3.5, -0.8)
  {$M_t / M_\infty$};

% Crank-style release curve: 1 - exp(-kt) approximation
\draw[engred, very thick] plot[domain=0:8.5, samples=40]
  (-3.5+\x, {-2.0 + 1.0 * (1 - exp(-0.4 * \x))});

\node[font=\scriptsize, anchor=west] at (-2.5, -1.3)
  {Cumulative release $\sim 1 - \exp(-k t)$ (long-time approx)};

\end{tikzpicture}
\caption[Drug-eluting stent cross-section]{Cross-section of a
drug-eluting stent strut with a polymer drug-reservoir coating of
thickness $L \approx 5\text{--}15\,\mu\text{m}$, in contact with
arterial tissue. The upper plot shows the radial drug
concentration $C(x, t)$ at an early time $t_1$ and at a late time
$t_2$: the coating depletes from a near-uniform initial loading,
and the tissue boundary layer carries the highest tissue
concentration. The lower plot shows the cumulative release
$M_t/M_\infty$ with a Crank-style one-dimensional diffusion solution;
the early-time regime is approximately Higuchi's
$\sqrt{t}$ law, the late-time regime is approximately
exponential with characteristic time $\tau \sim L^2/D$. The
design objective is the antiproliferative drug
window: too fast and the tissue concentration spikes above the
toxic threshold; too slow and the proliferative response is not
suppressed and the stent re-stenoses
\cite{paper:v6c12-crank-diffusion-1975,
paper:v6c12-higuchi-drug-release-1961,
paper:v6c12-stone-des-2010,
paper:v6c12-finn-des-pathology-2007}.}
\label{fig:vol06:ch12:des-cross-section}
\end{figure}
